% ! TeX root = ./main.tex

\section{Related Work}
\label{sec:amr:relwork}


\parahead{Variability and Runtime Analysis}
Performance variability in HPC environments is well documented~\cite{Petri2003:Missing,
  Hoefl2012:OptimizationPrinciplesCollective,Yildi2016:IONoise,Acun2016:Variation,VanS2009:AMR}, with
root causes ranging from tuning to contention.
Separately, fail-slow hardware has been reported as a source of performance
degradation~\cite{Gunaw2018:FailSlow,Lu2023:PERSEUS}.
Similar variability challenges have been identified in the context of collective performance in
hyperscale ML clusters~\cite{Gangi2024:RDMA}.
Varbench~\cite{Kocol2018:Varbench} is a framework to characterize computational variability using a
suite of compute kernels.
% Dedicated variability testbeds~\cite{Kocol2018:Varbench} and auto-tuning
% frameworks~\cite{Menon2020:AutotuningParameterChoices} have emerged to address these concerns.
Klenk et al.~\cite{Klenk2017:OverviewMPICharacteristics} analyze proxy MPI traces and report
significant MPI overheads (60\% at 1024 ranks) and attribute them to load imbalance, while also
noting limitations of trace analyses.
% observed rising MPI overheads at scale (doubling beyond 500 ranks) and attributed them to load
% imbalance.
While prior work documents separate causes of stragglers, an end-to-end treatment combining tuning
and placement optimization has been missing---a gap addressed by our work.
% However, such variability is typically not incorporated into application-level performance
% analysis or placement strategies. Our work addresses this gap by conducting an end-to-end study
% integrating telemetry collection, tuning, and placement.

\parahead{Performance Telemetry and Tooling}
Established performance tools such as TAU~\cite{Shend2006:TAU} and Score-P~\cite{Mey2012:Score-P}
collect profiles and traces (e.g., OTF2), but lack the queryability required for scalable
analytics, a limitation previously noted by Klenk et
al.~\cite{Klenk2017:OverviewMPICharacteristics}.
Caliper~\cite{Boehm2016:Caliper} supports structured telemetry and programmable aggregation, but
relies on plaintext data formats.
Parquet~\cite{:ApacheParquet} provides a compact, binary columnar format with embedded statistics
for efficient analytics, but is not widely adopted in performance tools.
Hatchet~\cite{Bhate2019:Hatchet} builds a structured in-memory representation from external
telemetry sources but assumes prior parsing.
Metric frameworks like LDMS~\cite{Schwa2021:LDMS} offer low-overhead system-level metrics, but
often lack application context needed for bottleneck attribution.
Yang et al.~\cite{Yang1988:CriticalPath} describe a general critical path analysis framework, while
our framework in \S\ref{sec:amr:tele:critp} is tailored for AMR codes.


\parahead{AMR and Communication Models}
% need to fix this
Modern AMR codes either use MPI or asynchronous runtimes like Charm++~\cite{Acun2014:Charm++} and
Legion~\cite{Bauer2012:Legion}.
% Boundary communication is often handled using point-to-point patterns, though n
Neighborhood collectives~\cite{Hoefl2011:ScalableProcessTopology,
  Hoefl2012:OptimizationPrinciplesCollective,Ghazi:EfficientDesignsDistributed} for boundary
communication have been proposed as an alternative to point-to-point messaging, but are currently
not used in the AMR codes we evaluated.
% These primitives, however, do not resolve placement or tuning challenges, especially under system
% variability~\cite{VanS2009:AMR}.

\parahead{Placement and Load Balancing}
Graph-based partitioners such as parMETIS~\cite{Lasal2013:parMETIS} and
Zoltan~\cite{Catal2007:Zoltan} handle complex meshes but incur high overhead.
Sasidharan et al.~\cite{Sasid2015:GeneralSpacefillingCurve} propose a SFC-based partitioner with
lower overhead than METIS.
All graph-based approaches model communication as edge cuts, which we find poorly correlated with
runtime communication overhead.
Meta-Balancer~\cite{Menon2012:MetaBalancer} focuses on rebalancing triggers, while Zheng et
al.~\cite{Zheng2011:PeriodicHierarchicalLoad} propose hierarchical schemes to scale load balancing.
Solver-based approaches are used to optimize placement in cloud datacenters~\cite{Newel2021:RAS},
but are computationally expensive (multi-hour) and impractical for AMR.



% Many established performance tools~\cite{Shend2006:TAU,Mey2012:Score-P} effectively capture
% profiles and event traces (OTF2)~\cite{Mey2012:Score-P}, but lack the queryability necessary for
% analytics, a limitation also noted in prior work~\cite{Klenk2017:OverviewMPICharacteristics}.
% Caliper~\cite{Boehm2016:Caliper} enables collecting structured and queryable telemetry with
% programmable aggregation, but uses plaintext formats like our ad hoc pipeline.
% Parquet~\cite{Apache:Parquet} is a highly efficient columnar format for analytics providing
% structured binary data with statistics and metadata, but is not widely used in performance
% tooling. Hatchet~\cite{Bhate2019:Hatchet} is an analysis library that parses structured telemetry
% into pandas dataframes augmented with a tree-based index for call-graphs, but needs to parse
% telemetry collected by other tools first. Separately, metric frameworks like
% LDMS~\cite{Schwa2021:LDMS} enhance queryability but focus on system-level state, often lacking
% application context for root-causing complex bottlenecks.

% Studies of performance variability include dedicated frameworks~\cite{Kocol2018:Varbench},
% analysis of causes like fail-slow hardware~\cite{Gunaw2018:FailSlow,Lu2023:PERSEUS}, and
% optimization of collectives~\cite{Hoefl2012:OptimizationPrinciplesCollective,
% Kanda2012:CanNetworkOffloadBased}. Related work explores proxy application
% limitations~\cite{Kruse2023:StatisticalAnalysisHPCa} and theoretical
% auto-tuning~\cite{Menon2020:AutotuningParameterChoices}.

% Modern AMR codes use MPI directly or task-based runtimes (e.g., Charm++~\cite{Acun2014:Charm++},
% Legion~\cite{Bauer2012:Legion}) for asynchronous execution. Neighborhood collectives have been
% proposed within MPI for communication patterns such as boundary communication~\cite{
% Hoefl2011:ScalableProcessTopology, Hoefl2012:OptimizationPrinciplesCollective,
% Ghazi:EfficientDesignsDistributed}, but do not inherently solve tuning challenges. % Key
% scalability bottlenecks include system variability and load balance~\cite{VanS2009:AMR}, with %
% strategies like Meta-Balancer~\cite{Menon2012:MetaBalancer} focusing on rebalancing triggers.

% Graph partitioning libraries (e.g., parMETIS~\cite{Lasal2013:parMETIS},
% Zoltan~\cite{Catal2007:Zoltan}) handle complex meshes but have higher overhead than the SFCs
% typical in AMR. Sasidharan et al.~\cite{Sasid2015:GeneralSpacefillingCurve} propose a SFC-based
% partitioning scheme with a lower overhead than Metis. Both approaches model meshes as graphs and
% use edge cuts as a proxy for communication load. We find the runtime impact of locality to be
% weakly correlated with theoretical models, and tradeable for better compute load balance. While
% hyperscaler placement solvers~\cite{Newel2021:RAS} exist, their multi-hour runtimes are unsuitable
% for the dynamism of AMR. Zheng et al.~\cite{Zheng2011:PeriodicHierarchicalLoad} propose
% hierarchical load-balancing schemes to circumvent scalability issues. Klenk et
% al.~\cite{Klenk2017:OverviewMPICharacteristics} analyze traces from proxy MPI applications,
% finding MPI overhead increases significantly with scale (\textgreater{}50\% at 500 ranks,
% \textgreater{}60\% at 1000 ranks), due to load imbalance. Our work demonstrates the value of
% integrating telemetry collection and analysis to understand and mitigate runtime artifacts.

% Variability in HPC environments is
% well-documented~\cite{Hoefl2012:OptimizationPrinciplesCollective,Yildi2016:IONoise,Petri2003:Missing,Acun2016:Variation}.
% Similar problems in modern ML clusters~\cite{Gangi2024:RDMA}. Van Straalen et
% al.~\cite{VanS2009:AMR} also identify system variability and load balance as key scalability
% bottlenecks, while approaches like Meta-Balancer~\cite{Menon2012:MetaBalancer} focus on when to
% trigger rebalancing.

% Established performance tools~\cite{Shend2006:TAU,Mey2012:Score-P} effectively capture profiles
% and event traces (OTF2)~\cite{Mey2012:Score-P}, but lack the inherent queryability found in
% analytical formats such as Parquet~\cite{Apache:Parquet}. Separately, structured metric collection
% frameworks like LDMS~\cite{Schwa2021:LDMS} enhance queryability but primarily focus on
% system-level metrics, and lack the fine-grained application context required for root-causing
% complex bottlenecks. %Substantial prior work exists on algorithms and techniques for parallel
% debugging~\cite{Yang1988:CriticalPath,Holli1998:CriticalPathProfiling,Aguil2003:Perf,Mitra2014:AccurateApplicationProgress}.
% Performance variability has been studied through frameworks like
% Varbench~\cite{Kocol2018:Varbench}, analysis of fail-slow
% hardware~\cite{Gunaw2018:FailSlow,Lu2023:PERSEUS}, and optimization of collective
% operations~\cite{Hoefl2012:OptimizationPrinciplesCollective, Kanda2012:CanNetworkOffloadBased}.
% Kruse et al.~\cite{Kruse2023:StatisticalAnalysisHPCa} study the limitations of proxy applications
% in helping tune real codes, while Menon et al.~\cite{Menon2020:AutotuningParameterChoices}
% describes a theoretical framework for automatic tuning. Modern AMR codes are built either directly
% on MPI or on task-based runtimes like Charm++~\cite{Acun2014:Charm++},
% Uintah~\cite{Holme2022:Uintah}, and Legion, which provide asynchronous execution abstractions.
% Neighborhood collectives have been proposed within MPI for communication patterns such as boundary
% communication~\cite{ Hoefl2011:ScalableProcessTopology,
% Hoefl2012:OptimizationPrinciplesCollective, Ghazi:EfficientDesignsDistributed}, but do not
% inherently solve tuning challenges. Van Straalen et al.~\cite{VanS2009:AMR} identify system
% variability and load balance as key scalability bottlenecks, while approaches like
% Meta-Balancer~\cite{Menon2012:MetaBalancer} focus on when to trigger rebalancing. General graph
% partitioning libraries like parMETIS~\cite{Lasal2013:parMETIS} and Zoltan~\cite{Catal2007:Zoltan}
% can handle complex mesh relationships but incur high overhead compared to space-filling curves.
% While hyperscalers use solvers to place workloads and services~\cite{Newel2021:RAS}, they take
% hours, making them unsuitable for AMR. Zheng et al.~\cite{Zheng2011:PeriodicHierarchicalLoad}
% propose hierarchical load-balancing schemes to circumvent scalability issues. Klenk et
% al.~\cite{Klenk2017:OverviewMPICharacteristics} analyze traces from proxy MPI applications,
% finding MPI overhead increases significantly with scale (\textgreater{}50\% at 500 ranks,
% \textgreater{}60\% at 1000 ranks), due to load imbalance. Our analysis demonstrates the value of
% colocating telemetry collection with analysis to isolate runtime artifacts.

% While hyperscale% placement systems~\cite{Newel2021:RAS} use mixed integer programming%
% complemented by temporary greedy heuristics, our experiments with commercial% solvers showed they
% cannot meet AMR timing constraints, taking hours versus our% 50ms target.%

% \textbf{Asynchronous Task-based Runtimes.} While some AMR codes are built directly on top of MPI,
% many AMR codes are built on top of frameworks that offer asynchronous task-based abstractions.
% Such frameworks allow users to model their application as dependency graphs of tasks, which are
% scheduled appropriately by the runtime. The goals of such frameworks are convenience, performance
% portability, overlapping communication and compution, and better separation of concerns. Examples
% of such frameworks are Charm++~\cite{Acun2014:Charm++}, Cello, Uintah~\cite{Holme2022:Uintah} and
% Legion. While it is possible for these frameworks to avoid common pitfalls, it is also possible
% for well-tuned codes directly using MPI to co-optimize the physics and the scheduling. No publicly
% available analyses and comparisons of the two approaches are available, to the best of our
% knowledge. \textbf{Graph stuff: } Zoltan~\cite{Catal2007:Zoltan}, METIS,
% ParMETIS~\cite{Lasal2013:parMETIS} RAS from Meta~\cite{Newel2021:RAS} We experimented with
% solver-based approaches but found them to be orders of magnitude slower than our target We tried
% an ILP formulation using a commercial state-of-the-art solver but failed to get solutions that
% outperformed simple heuristics even after 200s --- three orders of magnitude slower than our 50ms
% target per invocation. ILP approaches \textbf{Variability}: Varbench~\cite{Kocol2018:Varbench},
% Chombo~\cite{VanS2009:AMR} \textbf{Trace Paper}: Klenk et
% al.~\cite{Klenk2017:OverviewMPICharacteristics} analyze a set of traces from proxy MPI
% applications. While their analysis does not cover many AMR codes, they find the time spent by
% applications in MPI operations increases with scale (50\% for $\ge 500$ ranks, 60\% for $\ge 1000$
% ranks), which they ascribe to load imbalance. Our analyses suggest that while load imbalance in
% codes is a concern, jitter also creates stragglers which manifest as elevated MPI operation times.
% The difference between the two can be hard to establish when telemetry analysis is separated from
% infrastructure. We echo their concerns on the need for more detailed profiling, and the challenges
% it poses at scale. They suggest using asynchronous collectives over synchronous ones, but we find
% that the synchronicity is necessary to ensure correctness, and you can not drop-in replace things
% without a significant rehaul of the code. Variability from various sources and its imapct on
% runtime for bulk-synchronous parallel codes has been studied. Soem work also exists in exploiting
% it. - Many money quotes, need to be worked throughput the paper "The larger the scale the more
% time is spent in MPI, as applications with more than 500 ranks show an average MPI time of 57\%"
% "Use async collectives" "Overlap communication and computation"

% - Paper: Scalability Challenges for Massively Parallel AMR Straalen et al "Variability in system
% software was an equally important source of performance lsos" "finer-grained load-balancing
% needed" "improve scalability by one to two orders of magnitude"

% - Paper: Dubey et al, survey: achieving balanced load distribution is difficult because
% multiphysics, solvers etc. - Performnace and scaling have been major concerns for SAMR codes -
% Discusses heuristic-based approaches/

% and Zoltan~\cite{Catal2007:Zoltan} are used for placement in some context. They can handle
% arbitrary mesh topologies and complex relationships, their computational overhead makes simpler
% SFC-based approaches preferable for structured meshes. Approaches using mixed integer programming
% have been studied by hyperscalers for virtual machine allocation, but they take an hour to compute
% a new placement. Even in their designs, they are complemented by greedy heuristics for a temporary
% allocation while a better one is recomputed. For dynamic mesh codes these are impractical.

% % Critical paths

% Critical paths:~\cite{Yang1988:CriticalPath},~\cite{Aguil2003:Perf} %% Tuning It is commonly
% understood that networks require tuning. %% Noise Significant prior work exists on characterizing
% variability in hardware using benchmarks. Varbench~\cite{Kocol2018:Varbench} is a framework to
% characterize variability using compute kernels to elicit on-chip factors that create variability.
% While the existence of load-balancing and variability are both accepted, a study solving both
% these problems has not been attempted before. Fail-slow~\cite{Gunaw2018:FailSlow,Lu2023:PERSEUS}
% Meta-Balancer~\cite{Menon2012:MetaBalancer} is a study of when to invoke load-balancing. In
% frequently refining AMR codes, this decision is not necessary since a refinement automatically
% triggers a redistribution. In codes where no automatic trigger exists, this may be necessary. %%
% Neighborhood collectives %% Scalability challenges Straalen et al AMR
% scalability~\cite{VanS2009:AMR} %% nonblocking collectives/useless Nonblocking
% collectives~\cite{Kanda2012:CanNetworkOffloadBased} Neighborhood
% collectives~\cite{Hoefl2012:OptimizationPrinciplesCollective}
